% VI26-DETAILED-PROBLEM-01
\begin{problem}[Legacy AG5-P4: every affine open of a locally Noetherian scheme is Noetherian]\label{prob:vi26-p01}
Let \(X\) be a scheme.  Prove the legacy theorem that \(X\) is locally Noetherian if and only if every affine open \(U=\Spec A\subseteq X\) has \(A\) Noetherian.
\par\smallskip\begingroup\emergencystretch=2em\noindent\textit{Provenance.} Recovered numbered legacy Problem AG5-P4; previously routed by VI-SCOPE-\allowbreak 041 but absent from the ledger.\par\endgroup
\end{problem}

\ifdefined\FullProblemDossiers
\begin{solution}
The reverse implication is immediate: if every affine open has Noetherian coordinate ring, then any
affine open cover witnesses that \(X\) is locally Noetherian.

Assume conversely that \(X\) has an affine open cover \(V_i=\Spec B_i\) with every \(B_i\)
Noetherian.  Fix an arbitrary affine open \(U=\Spec A\subseteq X\).  For each \(x\in U\), choose
\(i\) with \(x\in V_i\).  Because distinguished opens form bases on both affine schemes, simultaneous
distinguished shrinking gives an open neighborhood
\[
W_x=D(f_x)\subseteq U,\qquad W_x=D(g_x)\subseteq V_i.
\]
Therefore
\[
A_{f_x}\cong (B_i)_{g_x},
\]
and the right-hand side is Noetherian because localization preserves Noetherianity.

The affine scheme \(U\) is quasi-compact, so finitely many of these neighborhoods cover:
\[
U=D(f_1)\cup\cdots\cup D(f_r).
\]
We now prove \(A\) Noetherian.  Let \(I\subseteq A\) be any ideal.  Since \(A_{f_j}\) is
Noetherian, \(I_{f_j}\) is finitely generated.  Choose finitely many elements of \(I\) whose
localizations generate \(I_{f_j}\) for each \(j\), and let \(J\subseteq I\) be the ideal generated
by the union of all these chosen elements.  Then
\[
I_{f_j}=J_{f_j}\quad\text{for every }j,
\]
so \(M:=I/J\) satisfies \(M_{f_j}=0\).

We make the annihilation step explicit.  For \(m\in M\), the equality \(M_{f_j}=0\) gives an exponent
\(N_j\) with \(f_j^{N_j}m=0\).  The distinguished opens \(D(f_j)\) cover \(\Spec A\), hence
\[
V(f_1,\ldots,f_r)=\varnothing
\quad\Longrightarrow\quad
(f_1,\ldots,f_r)=A.
\]
The same is true after replacing each \(f_j\) by a positive power, because the radicals are
unchanged; thus
\[
(f_1^{N_1},\ldots,f_r^{N_r})=A.
\]
Write \(1=\sum a_j f_j^{N_j}\).  Multiplying by \(m\) gives
\[
m=\sum a_j f_j^{N_j}m=0.
\]
Hence \(M=0\), so \(I=J\) is finitely generated.  Every ideal of \(A\) is therefore finitely
generated, and \(A\) is Noetherian.
\end{solution}
\fi
